\scanpage{35}
d'ambiguïté de signe), on trouve que pour un tel choix, disons
\[
\deleted{\unsure{\rho_0}}\quad\rho_0\mapsto\lambda_0,\quad\rho_1\mapsto\lambda_1,\quad\rho_\infty\mapsto\lambda_\infty\tag*{(87)}
\]
ou
\[
\tau_0\mapsto r_0,\quad\tau_1\unsure{\longrightarrow|\text{---}}r_1,\quad\tau_\infty\mapsto r_\infty,\tag*{(87 bis)}
\]
les autres choix possibles s'obtiennent en faisant des changements de signe sur les valeurs des images des $\rho_i$ resp. des $\tau_i$, avec comme seule contrainte que dans le cas des $\rho_i$, le nb de signes $-$ soit pair (donc égal à 2, \deleted{\unsure{3}}s'il y a effectivement modification). D'ailleurs, une fois le choix fait pour les $\rho_i$, si on veut que $\varphi^\circ$ soit induit par $\varphi$, cela impose le choix des images des $\tau_i$, au changement de signes \uline{simultané} près. Ceci vu, on tient à ce que l'image de $\rho_0$ soit \inserted{une matrice} unipotente, ce qui élimine l'ambiguïté de signe pour $\lambda_0$. \deleted{Cela} Il reste donc comme seule autre possibilité, en plus de (87), la possibilité
\[
\rho_0\mapsto\lambda_0,\ \rho_1\mapsto-\lambda_1,\ \rho_\infty\mapsto-\lambda_\infty\tag*{(88)}
\]
qui est celle qui s'est présentée dans (83). Elle se distingue du choix (87) par le fait que
\[
(-\lambda_\infty)^3=-(\lambda_\infty)^3=1,\quad\text{mais }\ (\lambda_\infty)^3=-1
\]
i.e. $-\lambda_\infty$ est d'ordre 3, tandis que $\lambda_\infty$ est d'ordre 6. Suivant le contexte, l'un ou l'autre des deux choix
% archive sheet number: 34
